% !PRINTED 517
\weberpage{517}
%<<CONT
sind. Bezeichnen wir also mit $p, p', p'', \ldots$ die Primfaktoren von $m$ von der Form $4n+1$, mit $q, q', q'', \ldots$ die Primfaktoren von $m$ von der Form $4n+3$, so kommen in den Teilgleichungen die folgenden Quadratwurzeln vor:
\[
\sqrt{p}, \sqrt{p'}, \sqrt{p''}, \ldots, \quad \sqrt{qq'}, \sqrt{qq''}, \ldots;
\]
wenn also nur ein $q$ in $m$ enthalten ist, so kommt $\sqrt{q}$ in den Teilgleichungen nicht vor. Da in diesem Fall wenigstens ein $q$ in $m$ aufgehen muß, so ist die Anzahl der zur Zerlegung erforderlichen Quadratwurzeln gleich der Anzahl der in $m$ aufgehenden Primzahlen, vermindert um $1$, also gleich der Anzahl der diesem Fall entsprechenden unabhängigen Charaktere.

2) Ist
\[
m\equiv 1\ (\mathrm{mod}\ 4), \quad m\equiv 6, 2\ (\mathrm{mod}\ 8), \quad m\equiv 4\ (\mathrm{mod}\ 16),
\]
so können nach § 105 beziehungsweise die Charaktere
\[
\left(\frac{-1}{A}\right), \ \left(\frac{2}{A}\right), \ \left(\frac{-2}{A}\right), \ \left(\frac{-1}{A}\right)
\]
durch die Charaktere von der Form
\[
\left(\frac{A}{m''}\right)
\]
ausgedrückt werden, und es reicht also auch in diesen Fällen die Formel (8) aus, um alle Geschlechter voneinander zu trennen.

Die Formeln (9), (12) lehren, daß in diesen Fällen zur vollständigen Trennung der Geschlechter folgende Adjunktionen nötig sind:
\weq{17}{
\begin{aligned}
m&\equiv 1\ (\mathrm{mod}\ 4), & \sqrt{p}, \sqrt{p'}, \sqrt{p''}, \ldots, \ & \sqrt{q}, \sqrt{q'}, \sqrt{q''}, \ldots \\
m&\equiv 6\ (\mathrm{mod}\ 8), & \sqrt{2}, \sqrt{p}, \sqrt{p'}, \sqrt{p''}, \ldots, \ & \sqrt{qq'}, \sqrt{qq''}, \ldots \\
m&\equiv 2\ (\mathrm{mod}\ 8), & \sqrt{p}, \sqrt{p'}, \sqrt{p''}, \ldots, \ & \sqrt{2q}, \sqrt{2q'}, \sqrt{2q''}, \ldots \\
m&\equiv 4\ (\mathrm{mod}\ 16), & \sqrt{p}, \sqrt{p'}, \sqrt{p''}, \ldots, \ & \sqrt{q}, \sqrt{q'}, \sqrt{q''}, \ldots
\end{aligned}
}

Denn im Fall $m\equiv 1 \ (\mathrm{mod}\ 4)$ ist die Anzahl der $q$ jedenfalls gerade. Setzt man also $m = q r^2 m''$, $m' = q r^2$ und versteht unter $r^2$ die größte in $r^2 m''$ aufgehende Quadratzahl, so ist $m''\equiv 3 \ (\mathrm{mod}\ 4)$ und die Formel (12) gibt die Adjunktion von $\sqrt{q}$, also ist in diesem Fall zur vollkommenen Trennung der Geschlechter
\[
\sqrt{p}, \sqrt{p'}, \sqrt{p''}, \ldots, \quad \sqrt{q}, \sqrt{q'}, \sqrt{q''}, \ldots
\]
zu adjungieren, und die Anzahl der Quadratwurzeln ist gleich der Anzahl der in $m$ aufgehenden Primzahlen, wieder in Über-
%CONT>>
